%% Contextual Learning Automata: Concept Outline
%% A working document to assess the idea for PhD contribution

\documentclass[11pt]{article}
\usepackage[utf8]{inputenc}
\usepackage[T1]{fontenc}
\usepackage[margin=1in]{geometry}
\usepackage{amsmath, amssymb, amsthm}
\usepackage{booktabs}
\usepackage{xcolor}
\usepackage{enumitem}
\usepackage{hyperref}
\usepackage{titlesec}

% Theorem environments
\newtheorem{definition}{Definition}
\newtheorem{proposition}{Proposition}
\newtheorem{question}{Open Question}

% Section formatting
\titleformat{\section}{\large\bfseries\color{blue!70!black}}{}{0em}{}[\titlerule]
\titleformat{\subsection}{\normalsize\bfseries}{}{0em}{}

\title{\textbf{Contextual Learning Automata}\\[0.5em]
\large Concept Outline and Novelty Assessment\\[0.3em]
\normalsize Working Document for PhD Contribution Evaluation}
\author{Ekaba Bisong}
\date{\today}

\begin{document}

\maketitle
\tableofcontents
\newpage

%% ============================================================
\section{The Problem: Why Context Matters}
%% ============================================================

\subsection{Your Current Approach: Global Policy LA}

In your current prototype, you have a \textbf{single Learning Automaton} that maintains one probability distribution $\mathbf{p} = [p_1, p_2, p_3, p_4]$ over four heuristics (NDVI, FAI, GNDVI, Ensemble).

\textbf{How it works:}
\begin{enumerate}
  \item For \textit{every} tile, sample an action $a \sim \mathbf{p}$
  \item Apply the selected heuristic, compute IoU reward $r$
  \item Update $\mathbf{p}$ based on $(a, r)$
\end{enumerate}

\textbf{The limitation:} The automaton learns a single ``best'' heuristic for the \textit{entire dataset}. It cannot adapt to individual tile characteristics.

\subsection{Why This Matters for Kelp Detection}

Different tiles have different characteristics:

\begin{center}
\begin{tabular}{lll}
\toprule
\textbf{Tile Feature} & \textbf{Condition} & \textbf{Likely Best Heuristic} \\
\midrule
Water turbidity & High (murky) & GNDVI (penetrates better) \\
Water turbidity & Low (clear) & FAI (designed for surface algae) \\
Cloud cover & Partial & Ensemble (robust to noise) \\
Kelp density & Dense canopy & NDVI (strong vegetation signal) \\
Kelp density & Sparse/submerged & FAI or GNDVI \\
\bottomrule
\end{tabular}
\end{center}

\textbf{Key insight:} The optimal heuristic depends on the \textit{context} (tile features). A global policy ignores this structure.

%% ============================================================
\section{Background: Two Separate Fields}
%% ============================================================

\subsection{Learning Automata (LA)}

Learning Automata are a classical approach from the 1970s-80s (Narendra \& Thathachar, 1989).

\begin{definition}[Learning Automaton]
A Learning Automaton is a tuple $(\mathcal{A}, \mathbf{p}, T)$ where:
\begin{itemize}
  \item $\mathcal{A} = \{a_1, \ldots, a_n\}$ is the action set
  \item $\mathbf{p} = [p_1, \ldots, p_n]$ is the probability vector over actions
  \item $T: \mathbf{p} \times \mathcal{A} \times \{0,1\} \to \mathbf{p}'$ is the update rule
\end{itemize}
\end{definition}

\textbf{Key property:} Standard LA is \textit{context-free}. The probability distribution $\mathbf{p}$ does not depend on any input features---it's the same for every decision.

\subsection{Contextual Bandits}

Contextual Bandits emerged from the machine learning community (Langford \& Zhang, 2007; Li et al., 2010).

\begin{definition}[Contextual Bandit]
At each round $t$:
\begin{enumerate}
  \item Observe context $\mathbf{x}_t \in \mathcal{X}$
  \item Select action $a_t \in \mathcal{A}$
  \item Receive reward $r_t$ (only for chosen action)
\end{enumerate}
The goal is to learn a policy $\pi: \mathcal{X} \to \mathcal{A}$ that maximizes cumulative reward.
\end{definition}

\textbf{Key property:} The action selection depends on the context $\mathbf{x}_t$.

\textbf{Common algorithms:} LinUCB, Thompson Sampling, $\varepsilon$-greedy with regression.

\subsection{The Gap Between These Fields}

\begin{center}
\begin{tabular}{lcc}
\toprule
\textbf{Property} & \textbf{Learning Automata} & \textbf{Contextual Bandits} \\
\midrule
Context-aware & No & Yes \\
Probability-based & Yes (core feature) & Sometimes \\
Convergence proofs & Strong (stationary env) & Regret bounds \\
Update mechanism & Simple, local & Often requires optimization \\
Computational cost & Very low & Can be high \\
Community & Control theory & Machine learning \\
\bottomrule
\end{tabular}
\end{center}

\textbf{Observation:} These two fields have developed largely independently. LA researchers don't typically consider context; bandit researchers don't typically use LA-style probability updates.

%% ============================================================
\section{Has This Been Done Before? Literature Review}
%% ============================================================

\subsection{Short Answer}

\textbf{Partially, but not in the way I'm proposing.}

\subsection{Related Work}

\textbf{1. Object Migration Automata (OMA)}

Oommen and others developed OMA for adaptive list organization. These use context in a limited sense (the object being accessed), but:
\begin{itemize}
  \item Context is discrete (object identity), not continuous features
  \item No general feature-based conditioning
  \item Different problem setting (list organization vs. action selection)
\end{itemize}

\textbf{2. Stochastic Point Location Automata}

Used for function optimization with continuous action spaces:
\begin{itemize}
  \item Actions are points in a continuous space
  \item But no input context---still learning a single optimal point
\end{itemize}

\textbf{3. Distributed/Networked Learning Automata}

Teams of automata that interact:
\begin{itemize}
  \item Each automaton is still context-free
  \item Coordination is via shared rewards, not via context
\end{itemize}

\textbf{4. Contextual Bandits with LA-style Updates}

To my knowledge: \textbf{No established work} directly combines:
\begin{itemize}
  \item Feature-based context conditioning
  \item LA probability update rules ($L_{R-I}$, $L_{R-P}$, Pursuit, etc.)
  \item Convergence analysis in the LA tradition
\end{itemize}

\subsection{Why Hasn't This Been Done?}

\begin{enumerate}
  \item \textbf{Community separation}: LA is from control theory; bandits are from ML. Different conferences, different terminology.
  \item \textbf{Different goals}: LA focuses on convergence to optimal action; bandits focus on regret minimization.
  \item \textbf{Computational assumptions}: LA assumes very limited computation (originated for hardware implementation); bandits assume modern computing.
\end{enumerate}

%% ============================================================
\section{Proposed Formulation: Contextual Learning Automata}
%% ============================================================

\subsection{Core Idea}

Instead of maintaining a single probability vector $\mathbf{p}$, maintain a \textbf{function} that maps context to probability:

\begin{definition}[Contextual Learning Automaton]
A Contextual Learning Automaton is a tuple $(\mathcal{A}, \mathcal{X}, \pi_\theta, T)$ where:
\begin{itemize}
  \item $\mathcal{A} = \{a_1, \ldots, a_n\}$ is the action set
  \item $\mathcal{X}$ is the context space (e.g., tile features)
  \item $\pi_\theta: \mathcal{X} \to \Delta^{n-1}$ is a parameterized policy mapping context to probability simplex
  \item $T$ is a parameter update rule inspired by LA
\end{itemize}
\end{definition}

\subsection{Concrete Instantiation: Linear Contextual LA}

Let context $\mathbf{x} \in \mathbb{R}^d$ be a feature vector for each tile.

\textbf{Policy:}
\[
\pi_\theta(\mathbf{x}) = \text{softmax}(\mathbf{W}\mathbf{x} + \mathbf{b})
\]
where $\theta = (\mathbf{W}, \mathbf{b})$ with $\mathbf{W} \in \mathbb{R}^{n \times d}$, $\mathbf{b} \in \mathbb{R}^n$.

\textbf{Action selection:}
\[
a_t \sim \pi_\theta(\mathbf{x}_t)
\]

\textbf{Update rule (LA-inspired):}
Instead of standard gradient descent, use LA-style updates:

For $L_{R-I}$ style:
\begin{itemize}
  \item If reward $r > 0$: update $\theta$ to increase $\pi_\theta(\mathbf{x}_t)[a_t]$
  \item If reward $r \leq 0$: no update (inaction)
\end{itemize}

Concretely, for the chosen action $a_t$ with context $\mathbf{x}_t$:
\[
\mathbf{W}_{a_t} \leftarrow \mathbf{W}_{a_t} + \alpha \cdot r \cdot \mathbf{x}_t \cdot (1 - \pi_\theta(\mathbf{x}_t)[a_t])
\]

\subsection{Alternative: Discretized Context}

For simpler analysis, discretize context into clusters:

\begin{enumerate}
  \item Cluster tiles into $K$ groups based on features (e.g., using k-means)
  \item Maintain $K$ separate probability vectors $\mathbf{p}^{(1)}, \ldots, \mathbf{p}^{(K)}$
  \item For tile with context $\mathbf{x}$, find cluster $k = \text{cluster}(\mathbf{x})$
  \item Use automaton with $\mathbf{p}^{(k)}$
\end{enumerate}

This is simpler to analyze theoretically (it's just $K$ independent automata) but loses the ability to generalize across similar contexts.

%% ============================================================
\section{What Would Be Novel?}
%% ============================================================

\subsection{Potential Theoretical Contributions}

\begin{enumerate}
  \item \textbf{Convergence analysis for Contextual LA}
  \begin{itemize}
    \item Under what conditions does the policy $\pi_\theta$ converge?
    \item What is the convergence rate?
    \item How does it compare to standard contextual bandits?
  \end{itemize}

  \item \textbf{Regret bounds with LA-style updates}
  \begin{itemize}
    \item Standard bandits use UCB or Thompson Sampling
    \item Can LA-style probability updates achieve sublinear regret?
    \item What are the trade-offs?
  \end{itemize}

  \item \textbf{Analysis of continuous rewards}
  \begin{itemize}
    \item Classical LA assumes binary rewards
    \item Your setting has continuous rewards (IoU $\in [0,1]$)
    \item How do convergence guarantees change?
  \end{itemize}

  \item \textbf{Non-stationary analysis}
  \begin{itemize}
    \item Kelp changes seasonally
    \item Can we prove adaptation rates under bounded drift?
  \end{itemize}
\end{enumerate}

\subsection{Potential Practical Contributions}

\begin{enumerate}
  \item \textbf{Demonstrate context improves performance}
  \begin{itemize}
    \item Show that Contextual LA outperforms Global LA on kelp detection
    \item Quantify the improvement
  \end{itemize}

  \item \textbf{Identify which context features matter}
  \begin{itemize}
    \item Learn which tile features predict optimal heuristic
    \item Interpretable model (unlike deep learning)
  \end{itemize}

  \item \textbf{Efficient online learning}
  \begin{itemize}
    \item LA updates are computationally simple
    \item Important for satellite imagery (large scale)
  \end{itemize}
\end{enumerate}

%% ============================================================
\section{Honest Assessment: Is This a PhD Contribution?}
%% ============================================================

\subsection{Arguments FOR}

\begin{enumerate}
  \item \textbf{Genuine gap in literature}: LA and contextual bandits have not been systematically combined.

  \item \textbf{Both theory and practice}: You can contribute both convergence analysis AND practical application to kelp detection.

  \item \textbf{Novel domain}: Applying LA to remote sensing heuristic selection is itself new.

  \item \textbf{Interpretability angle}: Unlike deep learning, LA provides interpretable probability dynamics.
\end{enumerate}

\subsection{Arguments AGAINST / Concerns}

\begin{enumerate}
  \item \textbf{``Just combining two things''}: A skeptical reviewer might say you're just putting context into LA, which is ``obvious.''

  \textit{Counter}: The theoretical analysis is non-trivial. LA convergence proofs rely on stationary environment and context-free updates. Extending to context requires new proof techniques.

  \item \textbf{Contextual bandits already exist}: Why not just use LinUCB or Thompson Sampling?

  \textit{Counter}: LA has different properties (probability-based, simple updates, specific convergence guarantees). The contribution is bringing these properties to the contextual setting.

  \item \textbf{Limited novelty if only empirical}: If you only show ``Contextual LA works better,'' that's a weak contribution.

  \textit{Counter}: You need theoretical analysis (convergence, regret bounds) to make it a strong contribution.
\end{enumerate}

\subsection{My Honest Opinion}

\textbf{This CAN be a PhD contribution, but it requires theoretical depth.}

A strong thesis would include:
\begin{enumerate}
  \item Formal definition of Contextual LA
  \item Convergence theorem (under stated assumptions)
  \item Regret analysis comparing to standard bandits
  \item Empirical validation on kelp detection
  \item Demonstration that context actually helps
\end{enumerate}

A weak thesis would be:
\begin{enumerate}
  \item ``I added context to LA''
  \item ``It works better on my dataset''
  \item No theoretical analysis
\end{enumerate}

%% ============================================================
\section{Alternative Directions (If Contextual LA Doesn't Work)}
%% ============================================================

If after investigation Contextual LA proves too similar to existing work, consider:

\subsection{Option A: Hierarchical LA}

Two-level automaton:
\begin{itemize}
  \item Level 1: Select detection \textit{scale} (e.g., 10m, 20m, 60m resolution)
  \item Level 2: Select heuristic at that scale
\end{itemize}

Novel: Theoretical analysis of hierarchical LA convergence.

\subsection{Option B: Non-Stationary LA with Drift Detection}

Focus on the temporal aspect:
\begin{itemize}
  \item Kelp appearance changes seasonally
  \item Design LA that detects distribution shift and adapts
\end{itemize}

Novel: Regret bounds under bounded non-stationarity.

\subsection{Option C: Spatially-Aware LA}

Focus on the spatial aspect:
\begin{itemize}
  \item Neighboring tiles share information
  \item Graph-regularized LA
\end{itemize}

Novel: Convergence analysis with spatial correlation.

%% ============================================================
\section{Concrete Next Steps}
%% ============================================================

If you want to pursue Contextual LA:

\begin{enumerate}
  \item \textbf{Literature deep-dive}: Search specifically for ``contextual learning automata,'' ``feature-dependent LA,'' ``state-dependent LA.'' Confirm the gap.

  \item \textbf{Implement discretized version}: Cluster tiles, run separate automata per cluster. See if context helps empirically.

  \item \textbf{Formalize the theory}: Write down the convergence theorem you want to prove. Identify the assumptions needed.

  \item \textbf{Discuss with advisor}: Get feedback on whether this direction is publishable and sufficient for PhD.
\end{enumerate}

%% ============================================================
\section{Open Questions for You to Consider}
%% ============================================================

\begin{question}
What tile features would you use as context? (Water indices, cloud mask values, geographic location, season?)
\end{question}

\begin{question}
How many distinct ``contexts'' exist in your data? Is discretization feasible?
\end{question}

\begin{question}
Are you comfortable with theoretical proof work, or do you prefer empirical focus?
\end{question}

\begin{question}
What does your advisor think about the LA direction generally?
\end{question}

\end{document}
